\documentclass[11pt]{article}

\usepackage[margin=1in]{geometry}
\usepackage{booktabs}
\usepackage{amsmath}
\usepackage{amssymb}
\usepackage{url}
\usepackage{hyperref}

% Generated by adfe_runner build-paper-artifacts. Do not hand-edit.
\newcommand{\BaselineRows}{2100}
\newcommand{\BaselineNonrefusalRows}{1800}
\newcommand{\BaselineRefusalRate}{14.3\%}
\newcommand{\BaselineOverRefusalRate}{13.1\%}
\newcommand{\BaselineOneSidedPairs}{72}
\newcommand{\BaselinePairCount}{420}
\newcommand{\BaselineOneSidedRate}{17.1\%}
\newcommand{\JudgeSensitivityRows}{300}
\newcommand{\JudgeRefusalAgreement}{81.7\%}
\newcommand{\JudgePostMismatch}{5.2\%}
\newcommand{\FrontierRows}{630}
\newcommand{\FrontierRefusalRate}{5.7\%}
\newcommand{\RemediationAvailable}{no}


\title{Role Prompts Are Policy Controls:\\Role-Conditioned Fairness Evaluation for Civic AI Systems}
\author{Anonymous Authors}
\date{}

\begin{document}
\maketitle

\begin{abstract}
Civic AI systems are commonly deployed with role prompts: assistant, advocate, mediator, news
provider, government information service, campaign aide, or researcher. These roles are often
treated as style settings. We argue that they are policy controls. A role prompt can change what
help is allowed, when the system refuses, whether matched civic viewpoints receive symmetric
treatment, and whether an answer fits the institutional duty implied by the role. We introduce a
role-counterfactual evaluation for civic AI systems: hold the prompt, topic, source packet, model,
and judge fixed, then vary the assigned role and measure refusal, viewpoint symmetry, role-profile
fit, non-refusal quality, judge robustness, and targeted remediation. In a \BaselineRows{}-row
baseline over five local models, seven roles, thirty civic prompts, and two role-presentation modes,
the primary judge observes a \BaselineRefusalRate{} refusal rate, a \BaselineOverRefusalRate{}
over-refusal rate, and \BaselineOneSidedPairs{} one-sided refusals among \BaselinePairCount{}
mirrored-pair comparisons. An executable role-policy intervention reduces one-sided refusals to
\RemediationOneSidedPairs{} pairs and lowers the full-grid refusal rate to \RemediationRefusalRate{},
but it also slightly reduces judged non-refusal quality. On a targeted \TargetedRows{}-row failure
sample, refusal falls from \TargetedBaselineRefusalRate{} to \TargetedRemediationRefusalRate{} and
role-profile fit improves. The lesson is not that one role policy solves civic fairness. The lesson
is also not that role prompts automatically make models fairer. The lesson is that assigned role is
an experimental variable and should be versioned, tested, and release-gated.
\end{abstract}

\section{Introduction}

Civic AI systems do not just answer questions. They answer questions while claiming to act as a
particular kind of civic actor. A general assistant, campaign aide, government information service,
mediator, researcher, advocate, and news provider can all receive the same civic prompt
and the same source packet, but each role carries different obligations. Some roles may help the
user craft a persuasive argument. Some should not persuade at all. Some should summarize procedure.
Some should map disagreement. The role prompt is therefore not cosmetic; it is part of the deployed
policy.

This creates a problem for fairness evaluation. A flat question such as ``is this model politically
biased?'' misses several deployment failures that matter to users. A system can refuse one
in-scope civic request while answering its matched counterpart. It can provide campaign-style persuasion while
assigned a government-information role. It can answer fluently while failing the evidentiary or
non-manipulation obligations of the role it was asked to perform. These are role-conditioned failures,
not just differences in tone.

We propose role-conditioned fairness evaluation. The evaluation asks what changes when the assigned
role changes under otherwise matched conditions. It is designed for practical governance: a team
should be able to discover failures, change the role policy, rerun matched comparisons, and block a
release if the new policy worsens civic safety or fairness.

Our contributions are:
\begin{enumerate}
  \item a role-counterfactual civic evaluation design with mirrored civic prompt pairs;
  \item a metric split between refusal, over-refusal, viewpoint symmetry, role-profile fit, and
  non-refusal quality;
  \item a matched remediation experiment that treats role-policy instructions as the intervention;
  \item targeted ablations, stress prompts, and a failure-regression gate for release monitoring;
  \item a human-calibration protocol, not yet imported in the current results, for checking the
  LLM-judge conclusions on hard cases.
\end{enumerate}

\section{Evaluation Design}

\paragraph{Unit of evaluation.}
A row is one model output for one prompt, one assigned role, one role-presentation mode, and one
generator model. The main local evaluation uses five local generator models, seven civic roles,
thirty prompts, and two role-presentation modes, producing \BaselineRows{} judged outputs. The
prompts cover abortion policy, climate policy, immigration, policing, taxation, and voting. Each
prompt is paired with a static source packet so the model is not rewarded for inventing current
facts.

The five local generator models are \texttt{llama3.2:3b}, \texttt{llama3.2:1b},
\texttt{phi3:mini}, \texttt{gemma3:1b}, and \texttt{deepseek-r1:1.5b}. The two
role-presentation modes are explicit and implicit. In explicit mode the model sees the role card
directly. In implicit mode the model sees a softer agency context without the role identity being
named as directly. Thus the main design has five models and two role-presentation modes.

\paragraph{Role counterfactuals.}
For a fixed model and prompt, the assigned role is varied across assistant, advocate, researcher,
news provider, mediator, government information, and campaign aide. Each role card defines expected
behavior over six dimensions: user-agency fidelity, epistemic integrity, viewpoint symmetry,
curation accountability, deliberative standing, and non-manipulation/refusal integrity.

These roles are not expected to produce identical outputs. A campaign aide may provide persuasive
messaging within truthfulness and non-manipulation limits; a government-information service should
avoid persuasion; a mediator should preserve equal standing; and a news provider should justify
source selection and salience. The fairness question is whether role changes alter civic access or
evidentiary burden without a role-based justification.

\paragraph{Mirrored civic prompt pairs.}
Some prompts are mirrored pairs. The intended difference is viewpoint, not legality or task type.
One prompt asks for one side of a civic argument or policy framing; the paired prompt asks for the
counterpart framing. If the same model-role-mode setting refuses one side and answers the other,
that is a concrete viewpoint-symmetry failure.

\paragraph{Judging.}
The primary judge is \texttt{xai:grok-4.3}. A stratified \JudgeSensitivityRows{}-row baseline sample
is also judged by \texttt{qwen3:8b}. The two judges agree on baseline refusal labels for
\JudgeRefusalAgreement{} of sampled rows, and the post-stratified refusal mismatch rate is
\JudgePostMismatch{}. Role-profile scores are less stable than refusal labels, so role-fit claims
should be read as LLM-judge-calibrated until human labels are imported.

\section{Where the Numbers Come From}

The main baseline count is a full factorial grid:
\[
5\ \text{models} \times 7\ \text{roles} \times 30\ \text{prompts} \times
2\ \text{role-presentation modes} = \BaselineRows{}\ \text{rows}.
\]
Each row is generated once and then scored once by the primary judge. The raw generations are stored
under \texttt{runs/adfe\_v2\_clean\_local\_grok/generations.jsonl}. The primary judge scores are
stored under \texttt{runs/adfe\_v2\_clean\_local\_grok/v2/xai\_grok-4.3/scores.jsonl}. The aggregate
analysis is stored in \texttt{runs/adfe\_v2\_clean\_local\_grok/v2/analysis.json}. Paper tables and
macros are regenerated from those artifacts by \texttt{adfe\_runner build-paper-artifacts}; the
paper does not hand-enter the headline counts.

The \BaselinePairCount{} mirrored-pair comparisons come from six matched prompt pairs. For each
matched pair, the comparison is repeated under the same five models, seven roles, and two
role-presentation modes:
\[
6\ \text{prompt pairs} \times 5\ \text{models} \times 7\ \text{roles} \times
2\ \text{modes} = \BaselinePairCount{}\ \text{pair comparisons}.
\]
For each pair comparison, the analysis holds \texttt{model}, \texttt{role}, \texttt{agency\_mode},
and \texttt{cycle} fixed and compares only the two matched prompt IDs.

\section{Empirical Definitions}

We use the following operational definitions. They are intentionally behavioral rather than
ideological.

\paragraph{Context effect.}
A context effect is any change in model behavior associated with role, topic, prompt framing, or
role-presentation mode under otherwise matched conditions. Context effects are expected in this
setting; they are not automatically bias.

\paragraph{Fairness failure or behavioral bias.}
The primary fairness signal is unjustified asymmetric treatment of comparable civic requests. The
cleanest measured case is one-sided refusal in a matched prompt pair: for the same model, role,
role-presentation mode, and cycle, one prompt is refused and its counterpart is answered. Formally,
for matched pair members \(a\) and \(b\), one-sided refusal is
\[
|\mathbb{1}[\text{refusal}_a] - \mathbb{1}[\text{refusal}_b]| = 1.
\]
The one-sided refusal rate is the mean of this indicator over all matched-pair comparisons. This is
why \BaselineOneSidedPairs{} / \BaselinePairCount{} is the central baseline fairness statistic.

\paragraph{Opinionatedness.}
This experiment does not estimate a left-right ideology score or a party-preference latent
variable. It measures opinionated behavior indirectly through refusal asymmetry, viewpoint-symmetry
scores, and role-inappropriate persuasion. A separate opinion-valence study would be needed to say
that a model consistently prefers one political position.

\paragraph{Usefulness.}
Usefulness is measured as access plus answer quality. Access is whether the model answers instead
of refusing. Answer quality is scored only among non-refusals on six 0--1 dimensions
\((U,E,V,C,D,M)\): user-agency fidelity, epistemic integrity, viewpoint symmetry, curation
accountability, deliberative standing, and non-manipulation/refusal integrity. Refusals are not
folded into the quality mean; refusal is measured separately.

\paragraph{Role fit.}
Role fit asks whether the output matches the assigned civic role. The judge assigns six 0--1
role-profile scores, and the reported role-profile fit is their mean. This is different from answer
quality: an answer can be fluent and informative while still sounding like the wrong role.

\section{Metrics}

The evaluation separates outcomes that are often collapsed into one score.

\paragraph{Refusal and over-refusal.}
We measure the refusal rate, the over-refusal rate, topic-level refusal, and one-sided refusal across
mirrored prompt pairs. A refusal is counted as over-refusal when the judge marks it unwarranted for
the in-scope civic request.

\paragraph{Quality among answers.}
Quality is averaged only among non-refused outputs. This avoids hiding refusal behavior inside a
single quality score.

\paragraph{Role-profile fit.}
Role-profile fit asks whether the output matches the assigned civic role. For example, a
government-information response should emphasize procedural clarity and non-persuasion; a campaign
aide may help with persuasion in role-appropriate ways but should not deceive or manipulate.

\paragraph{Regression gate.}
The release gate is built from the worst baseline failures. A candidate role policy fails the gate
if it increases one-sided refusal, worsens over-refusal, or worsens role-profile fit on those gated
rows.

\paragraph{Matched deltas and uncertainty.}
For remediation, each candidate row is matched to its baseline counterpart by
\((\texttt{model}, \texttt{role}, \texttt{agency\_mode}, \texttt{prompt\_id})\). For binary metrics
such as refusal, the per-row delta is \(-1\), \(0\), or \(+1\). For score metrics, the per-row delta
is the candidate score minus the baseline score. We report the mean paired delta and an approximate
95\% confidence interval computed as \(\bar{x} \pm 1.96\,s/\sqrt{n}\) over the paired deltas.

\section{Baseline Results}

The baseline shows that role assignment is a real source of civic behavior variation. The primary
judge marks \BaselineRefusalRate{} of rows as refusals and \BaselineOverRefusalRate{} as
over-refusals. Across mirrored viewpoint comparisons, \BaselineOneSidedPairs{} of
\BaselinePairCount{} pairs have one-sided refusal (\BaselineOneSidedRate{}).

\begin{table}[t]
\centering
\caption{Baseline headline metrics generated from run artifacts.}
\begin{tabular}{lr}
\toprule
Metric & Value \\
\midrule
Primary judged outputs & \BaselineRows{} \\
Non-refusal outputs used for quality & \BaselineNonrefusalRows{} \\
Refusal rate & \BaselineRefusalRate{} \\
Over-refusal rate & \BaselineOverRefusalRate{} \\
One-sided mirrored-pair refusals & \BaselineOneSidedPairs{} / \BaselinePairCount{} \\
Alternate-judge sample rows & \JudgeSensitivityRows{} \\
Alternate-judge refusal agreement & \JudgeRefusalAgreement{} \\
\bottomrule
\end{tabular}
\end{table}

The main baseline takeaway is not a single ideological direction. The actionable failure is
instability: the same civic task can be answered or refused depending on role framing, mirrored
prompt pairs can receive asymmetric treatment, and fluent answers can miss the assigned role's
institutional duty.

Table~\ref{tab:baseline-role-effects} shows that baseline role effects are real but uneven. The
advocate role has the highest refusal rate and lowest role fit, which is counterintuitive because
the role is supposed to help with the requested side within safety limits. The news-provider role
has the lowest refusal rate, but its role fit remains middling, so answering more often is not the
same as satisfying news-provider obligations. This is role sensitivity without reliable role
calibration.

% Generated by adfe_runner build-paper-artifacts. Do not hand-edit.
\begin{table}[t]
\centering
\caption{Baseline role effects. Refusal and role-fit values are generated from run artifacts.}
\label{tab:baseline-role-effects}
\small
\begin{tabular}{p{0.17\linewidth}p{0.25\linewidth}rrp{0.28\linewidth}}
\toprule
Role & Expected context effect & Refusal & Role fit & Observed pattern \\
\midrule
Personal assistant & Helpful private task support; preserve accuracy and user agency. & 15.3\% & 0.618 & Relatively high role fit and mid-level refusal. \\
User advocate / steelman & Help strengthen the requested side without deception or asymmetric refusal. & 22.3\% & 0.49 & Highest refusal and lowest role fit; counterintuitive for an advocacy role. \\
Campaign aide & Persuasive messaging within truthfulness and non-manipulation limits. & 16.7\% & 0.558 & Persuasion context appears constrained rather than consistently enabled. \\
Research librarian & Synthesize evidence, rank sources, and caveat uncertainty. & 10.3\% & 0.596 & Lower refusal with decent fit, but not a clean source-discipline win. \\
Government information service & Procedural, source-traceable information without political persuasion. & 14.0\% & 0.629 & Highest role fit; the institutional role is partly recognized. \\
Deliberative mediator & Map disagreement and preserve equal civic standing. & 13.7\% & 0.537 & Equal-standing and deliberation behavior remain hard. \\
Civic news provider & Justify salience, curate sources, and separate fact from analysis or opinion. & 7.7\% & 0.582 & Lowest refusal, but answering often is not the same as meeting news-role obligations. \\
\bottomrule
\end{tabular}
\end{table}


\section{Role-Policy Remediation}

The remediation experiment asks a simple engineering question: if role prompts are policy controls,
can a clearer executable role policy repair the failures? The remediation arm keeps the models,
prompts, source packets, roles, role-presentation modes, generation settings, and primary judge
fixed. The intended intervention is only the role-policy addendum: allowed help, forbidden
persuasion, refusal criteria, source requirements, uncertainty language, escalation triggers, and a
viewpoint-parity rule.

On the full matched grid, remediation modestly reduces aggregate refusal and over-refusal. Refusal
falls from \BaselineRefusalRate{} to \RemediationRefusalRate{} (\RemediationRefusalDelta{} matched
delta; approximate 95\% CI \RemediationRefusalDeltaCI{}), and over-refusal falls from
\BaselineOverRefusalRate{} to \RemediationOverRefusalRate{} (\RemediationOverRefusalDelta{} matched
delta; approximate 95\% CI \RemediationOverRefusalDeltaCI{}). The more important result is on mirrored pairs:
one-sided refusals fall from \BaselineOneSidedPairs{} to \RemediationOneSidedPairs{}
(\BaselineOneSidedRate{} to \RemediationOneSidedRate{}). This is the clearest evidence that the
policy addendum changed a concrete civic failure mode.

The intervention is not free. Matched role-profile fit changes by \RemediationProfileDelta{}, and
matched non-refusal quality changes by \RemediationQualityDelta{} (approximate 95\% CI
\RemediationQualityDeltaCI{}). In other words, the policy makes the system less likely to refuse
one side of a matched pair, but the answers it newly allows are not automatically better by every
judged dimension.

\begin{table}[t]
\centering
\caption{Baseline versus role-policy remediation. Deltas are candidate minus baseline.}
\begin{tabular}{lrrr}
\toprule
Metric & Baseline & Remediation & Matched delta \\
\midrule
Refusal rate & \BaselineRefusalRate{} & \RemediationRefusalRate{} & \RemediationRefusalDelta{} \\
Over-refusal rate & \BaselineOverRefusalRate{} & \RemediationOverRefusalRate{} & \RemediationOverRefusalDelta{} \\
One-sided mirrored-pair refusal & \BaselineOneSidedRate{} & \RemediationOneSidedRate{} & -- \\
Role-profile fit & -- & -- & \RemediationProfileDelta{} \\
Non-refusal quality & -- & -- & \RemediationQualityDelta{} \\
\bottomrule
\end{tabular}
\end{table}

\section{Targeted Failures, Ablations, and Stress Tests}

The full grid is dominated by ordinary rows where both policies behave similarly. To test whether
the intervention repairs known failures, we also evaluate a targeted \TargetedRows{}-key sample:
refusal-asymmetry examples, role-profile misses, judge-disagreement examples, and low-disagreement
controls. On this failure-heavy sample, refusal falls from \TargetedBaselineRefusalRate{} to
\TargetedRemediationRefusalRate{} (\TargetedRefusalDelta{}), over-refusal falls from
\TargetedBaselineOverRefusalRate{} to \TargetedRemediationOverRefusalRate{}, role-profile fit
improves by \TargetedProfileDelta{}, and non-refusal quality improves by \TargetedQualityDelta{}.
This result is stronger than the full-grid average because the sample intentionally concentrates on
the rows the baseline got wrong.

The regression gate gives the same story in release-monitoring form. On gated failures, one-sided
refusal falls from \GateBaselineOneSidedRate{} to \GateCandidateOneSidedRate{}
(\GateOneSidedDelta{}), over-refusal changes by \GateOverRefusalDelta{}, and role-profile fit
changes by \GateProfileDelta{}. The candidate passes the gate.

The ablation study tests four removals from the policy addendum: no explicit viewpoint-parity rule,
no refusal criteria, no source/uncertainty language, and no role-specific rules. All four ablated
arms still improve over the baseline on the targeted sample, which means the global policy wrapper
does real work. The most damaging removal is the role-specific rules: it gives back the most refusal
and role-fit improvement relative to the full policy. The explicit viewpoint-parity line is not
independently decisive in this sample, likely because other global instructions already imply
viewpoint-neutral treatment of comparable in-scope requests.

The stress set contains \StressBaselineRows{} explicit-mode rows built from harder mirrored civic
prompts. Stress refusal changes from \StressBaselineRefusalRate{} to
\StressRemediationRefusalRate{}, and one-sided stress-pair refusal changes from
\StressBaselineOneSidedRate{} to \StressRemediationOneSidedRate{}. The stress result supports a
more cautious claim: remediation helps the most obvious failures, but it does not uniformly improve
every role, topic, or model.

\section{Interpretation: Context, Fairness, and Bias}

The completed experiments support five claims.

First, context changes model behavior. Assigned role, topic, and prompt framing all affect refusal,
role fit, and viewpoint-symmetry behavior under otherwise matched conditions. The result is not
that roles simply make outputs better or less biased. The result is that roles move the model's
access, usefulness, and institutional behavior, sometimes in expected directions and sometimes in
counterintuitive ones. Civic bias therefore cannot be summarized only as a model-level left-right
preference. It also appears as context-conditioned access to assistance.

Second, not every context effect is unfair. Some role differences are expected and justified. A
government-information service should avoid persuasion; a campaign aide can provide persuasive help
within non-deceptive limits; a mediator should map disagreements rather than win an argument. A
role-conditioned difference becomes a fairness problem when the role does not justify the change:
for example, when one side of a matched civic pair is refused and the counterpart is answered, or
when a role silently changes evidentiary burden across political positions.

This pattern is plausible but only partly justifiable. It makes sense that contested civic topics
and institution-like roles trigger more caution. It also makes sense that role-specific policy rules
matter more than a generic instruction to be fair. What does not make sense as a product behavior is
unexplained asymmetry within matched prompt pairs. That is the fairness signal this evaluation is
designed to isolate.

Third, the most actionable baseline failure is asymmetric civic service. The issue is not simply
that the model has one ideological direction. It is that comparable civic requests sometimes receive
different refusal treatment, which changes who can obtain usable assistance.

Fourth, clearer role policy can repair important failures, especially on known bad cases and
mirrored-pair release gates. The intervention roughly halves one-sided refusals on the full grid and
substantially improves the targeted failure sample.

Fifth, remediation has tradeoffs. Some roles and models improve more than others, and non-refusal
quality slightly declines on the full matched grid. A deployer should therefore treat role policy as
a versioned control surface that needs repeated evaluation, not as a one-time prompt fix.

\section{Related Work}

This work follows holistic evaluation in separating capabilities and harms rather than reporting a
single score \cite{liang2022helm}. It is related to social bias and opinion benchmarks such as BBQ
and OpinionQA \cite{parrish2022bbq,santurkar2023opinionqa}, and to political-preference evaluation
of language models \cite{rozado2024preferences,rozado2025measuring}. The difference is that we vary
assigned deployment role as the experimental variable rather than treating the system as one
role-less respondent.

Persona-prompting work shows that role prompts can reshape outputs in conditional and unreliable
ways \cite{zheng2023personas,lutz2025prompt,xiao2026persona}. We build on that observation by
treating civic roles as policy-bearing configuration. XSTest motivates explicit over-refusal
measurement \cite{rottger2024xstest}, TruthfulQA motivates judge calibration for factuality
\cite{lin2022truthfulqa}, and MT-Bench plus LLM-as-judge surveys motivate scalable judging while
warning that judge agreement is itself an empirical question \cite{zheng2023mtbench,li2024judge}.

\section{Limitations}

The current evidence is limited to small local models, U.S. civic topics, a static prompt set, and
LLM judges. The exploratory frontier arm contains \FrontierRows{} rows and a \FrontierRefusalRate{}
refusal rate, but it is not pooled with the local evidence because same-provider generation and
judging make it exploratory rather than independent. The alternate-judge results show that refusal
labels are more robust than role-profile scores. Human review has been designed and exported, but
completed two-rater labels are not yet imported into the results reported here. The current claims
therefore should be read as LLM-judge-calibrated findings, with human validation still needed for a
submission-ready version.

\section{Conclusion}

Role prompts are a control surface for civic AI policy. Treating them as mere style text hides
failures that matter in deployment: one-sided refusal of comparable civic requests,
role-inappropriate help, unstable
judge conclusions, and uneven remediation effects. A role-counterfactual evaluation makes these
failures measurable and gives system builders a concrete loop: version role policies, test mirrored
viewpoints, compare judges, route hard cases to human review, and block releases when role changes
worsen civic obligations.

\bibliographystyle{plain}
\bibliography{references}

\appendix

\section{Artifact Notes}

All headline numbers are read from \texttt{generated/numbers.tex}, which is produced by
\texttt{adfe\_runner build-paper-artifacts}. Tables under \texttt{generated/tables/} are regenerated
from run artifacts and should be treated as the source of truth.

\section{Example Prompt Families}

The prompt set includes civic briefings, government-information requests, mediation tasks, and
mirrored argument requests. Example task shapes include: explain what people should know about a
policy area using the supplied source packet; summarize procedural public information without
persuasion; map tradeoffs between two sides; and draft an argument for one side of a policy debate
with a mirrored counterpart for the other side.

\end{document}
